% from ega4-8-fr.tex
\oldpage[IV-3]{39}
En vertu de (8.6.3), il existe donc un $\lambda\geq\alpha$ et un
sous-préschéma (resp. induit sur un ouvert, resp. fermé) $X'_\lambda$ de
$S_\lambda$, de présentation finie sur $S_\lambda$, tels que $X'$ soit
isomorphe à $X'_\lambda\times_{S_\lambda}S$. Pour tout $\mu\geq\lambda$,
$X'_\mu=X'_\lambda\times_{S_\lambda}S_\mu$ est donc un sous-préschéma (resp.
induit sur un ouvert, resp. fermé) de $S_\mu$, de présentation finie sur
$S_\mu$, et il résulte donc de (8.8.2.4) et (8.8.2.5) qu'il existe un
$\mu\geq\lambda$ et un isomorphisme $g_\mu:X_\mu\to X'_\mu$ tel que
$g=g_\mu\times1_S$ soit l'isomorphisme $X\to X'$ associé à $f$~; d'où la
conclusion.

\emph{Cas (vi) et (vii)}~: On sait (1.8.4.1) que
$Z_\alpha=f_\alpha(X_\alpha)$ est constructible dans $S_\alpha$~; si l'on
pose $Z_\lambda=u_{\alpha\lambda}^{-1}(Z_\alpha)$ pour $\lambda\geq\alpha$
et $Z=u_\alpha^{-1}(Z_\alpha)$, on a $Z_\lambda=f_\lambda(X_\lambda)$ et
$Z=f(X)$ (I, 3.4.8). Comme en vertu de (8.3.11) l'application canonique
$\varinjlim\mathfrak C(S_\lambda)\to\mathfrak C(S)$ est injective, la
relation $f(X)=S$ implique l'existence d'un $\lambda\geq\alpha$ tel que
$f_\lambda(X_\lambda)=S_\lambda$, ce qui prouve le théorème dans le cas (vi).
Pour le prouver dans le cas (vii), il suffit de noter que le morphisme
structural $X_\alpha\times_{S_\alpha}X_\alpha\to S_\alpha$ est de présentation
finie puisqu'il en est ainsi de la première projection
$X_\alpha\times_{S_\alpha}X_\alpha\to X_\alpha$ (1.6.2)~; il suffit donc en
vertu de (1.8.7.1) d'appliquer le cas (vi) du théorème au morphisme diagonal
$\Delta_{f_\alpha}:X_\alpha\to X_\alpha\times_{S_\alpha}X_\alpha$ en notant
que l'on a $\Delta_{f_\lambda}=\Delta_{f_\alpha}\times1_{S_\lambda}$ et
$\Delta_f=\Delta_{f_\alpha}\times1_S$ (I, 5.3.4 et 3.3.11).

\emph{Cas (viii) et (ix)}~: Comme $S=\operatorname{Spec}(A)$ est affine, dire
que $f:X\to S$ est affine (resp. quasi-affine) signifie qu'il existe un entier
$r$ et une immersion fermée (resp. une immersion)
$j:X\to\mathbf V_S^r=\operatorname{Spec}(A[T_1,\ldots,T_r])$ de
$S$-préschémas, puisque $f$ est de type fini et $S$ quasi-compact
(II, 5.1.9). Comme
$\mathbf V_S^r=\mathbf V_{S_0}^r\times_{S_0}S$, et que
$\mathbf V_{S_0}^r$ est un $S_0$-préschéma de présentation finie, il résulte
de (8.8.2, (i)) appliqué au $S$-morphisme $j$, qu'il existe un $\lambda$ et un
$S_\lambda$-morphisme $j_\lambda:X_\lambda\to\mathbf V_{S_\lambda}^r$ tel que
$j=j_\lambda\times1_S$~; appliquant alors (iv) (resp. (ii)) à $j$, on en
déduit qu'il existe $\mu\geq\lambda$ tel que $j_\mu$ soit une immersion
fermée (resp. une immersion)~; par suite $f_\mu$ est affine (resp.
quasi-affine).

\emph{Cas (x)}~: Par hypothèse, on a $X=\operatorname{Spec}(B)$, où $B$ est
une $A$-algèbre qui est un $A$-module de présentation finie (1.4.7)~; il
résulte donc de (8.5.2, (ii)) qu'il y a un $\lambda$ et un $A_\lambda$-module
de présentation finie $B_\lambda$ tels que le $A$-module $B$ soit isomorphe à
$B_\lambda\otimes_{A_\lambda}A$. La structure de $A$-algèbre de $B$ est
définie par un $A$-homomorphisme $m:B\otimes_AB\to B$~; comme on a
$B\otimes_AB=(B_\lambda\otimes_{A_\lambda}B_\lambda)\otimes_AA$, il existe
d'après (8.5.2, (i)) un $\mu\geq\lambda$ et un $A_\mu$-homomorphisme
$m_\mu:B_\mu\otimes_{A_\mu}B_\mu\to B_\mu$ tel que $m=m_\mu\otimes1$. En
considérant les diagrammes usuels exprimant l'associativité et la
commutativité de $m$, on voit en appliquant encore (8.5.2, (i)) qu'il existe
$\nu\geq\mu$ tel que $m_\nu$ définisse sur $B_\nu$ une multiplication
associative et commutative~; de la même manière on voit qu'on peut supposer
$\nu$ pris assez grand pour que l'anneau $B_\nu$ ainsi défini admette un
élément unité. Si $X'_\nu=\operatorname{Spec}(B_\nu)$, il est clair alors que
$X$ est $S$-isomorphe à $X'_\nu\times_{S_\nu}S$, donc, en vertu de (i), il
existe $\rho\geq\nu$ tel que $X_\rho$ et $X'_\rho$ soient
$S_\rho$-isomorphes, ce qui achève la démonstration dans ce cas.

Pour démontrer le théorème dans le cas (xii), nous démontrerons d'abord la
proposition suivante~:
\end{proof}

\begin{proposition}[8.10.5.1]
\label{IV.8.10.5.1-fr}
(lemme de Chow pour les morphismes de présentation finie). --- Soient $A$ un
anneau, $X$, $Y$ deux $A$-préschémas de présentation finie,
$f:X\to Y$
